\section{Task 3 -- Thiết kế kiến trúc}\label{ch:kientruc}

Tạo phẩm của Task 3 nằm trong \code{docs/\allowbreak{}architecture/\allowbreak{}}:
\begin{itemize}
\item \code{architecture.\allowbreak{}md}: mô tả kiến trúc;
\item \code{container-diagram.\allowbreak{}puml}: sơ đồ C4 mức container;
\item \code{deployment-diagram.\allowbreak{}puml}: deployment diagram;
\item 8 ADR trong \code{adr/}.
\end{itemize}

Task 3 được merge qua PR \#8 sáng 15/11/2026, ngay trước khi khởi tạo mã nguồn ở Task 4 (PR \#9).

\subsection{Task 3.1 -- Phương pháp và kiến trúc tổng thể}
Nhóm dùng \textbf{kiến trúc phân lớp} kết hợp \textbf{chia module theo domain} ở backend, còn gọi là modular monolith, và trình bày theo hai mức của mô hình C4:
\begin{itemize}
\item \textbf{Mức Container:} trình duyệt, ứng dụng Next.js, API NestJS, PostgreSQL, cổng thanh toán.
\item \textbf{Mức Component:} các module NestJS bên trong API.
\end{itemize}

Đầu vào của thiết kế là SRS, UC-01 và các sơ đồ của Task 2. Đặc biệt, class diagram quyết định ranh giới module.

\hinh{container-diagram.pdf}{\linewidth}{Sơ đồ C4 mức container (\code{docs/\allowbreak{}architecture/\allowbreak{}container-diagram.\allowbreak{}puml})}{fig:container}

{\small
\begin{longtable}{@{}L{2.2cm}L{2.6cm}L{4.6cm}L{2.4cm}L{2.1cm}@{}}
\caption{Các container của hệ thống}\label{tab:container}\\
\toprule
\textbf{Container} & \textbf{Công nghệ} & \textbf{Trách nhiệm} & \textbf{Mã nguồn} & \textbf{Yêu cầu chính} \\ \midrule \endfirsthead
\toprule \textbf{Container} & \textbf{Công nghệ} & \textbf{Trách nhiệm} & \textbf{Mã nguồn} & \textbf{Yêu cầu chính} \\ \midrule \endhead
\bottomrule \endfoot
Web app & Next.js (App Router), React, TypeScript & Giao diện mobile-first: menu, giỏ hàng, đặt đơn, thanh toán, lịch sử, màn hình nhân viên. Giỏ hàng lưu ở trình duyệt & \code{apps/web} & FR-01..04, FR-11, FR-15, NFR-05, NFR-06 \\
API & NestJS, TypeScript, Prisma & Toàn bộ nghiệp vụ: xác thực, đơn hàng, thanh toán, tồn kho, khuyến mãi, điểm thưởng. REST/JSON & \code{apps/api} & FR-05..FR-21, NFR-02, NFR-04 \\
Database & PostgreSQL 16 & Lưu dữ liệu theo ERD; transaction cho đặt đơn và thanh toán & \code{apps/\allowbreak{}api/\allowbreak{}prisma} (schema, migration, seed) & FR-12, NFR-10 \\
Cổng thanh toán mô phỏng & Module trong API (\code{MockPaymentGateway}) & Xử lý Ví, Thẻ, QR với tài khoản thử nghiệm; điều khiển được kết quả để kiểm thử & \code{apps/\allowbreak{}api/\allowbreak{}src/\allowbreak{}payments} & FR-08, FR-18 \\
VNPay sandbox & Dịch vụ bên ngoài, HTTPS & Thanh toán thử nghiệm qua redirect và Return URL & Adapter \code{VnpayGateway} trong \code{apps/\allowbreak{}api/\allowbreak{}src/\allowbreak{}payments} & FR-21 (Could) \\
\end{longtable}
}

Luồng gọi:
\begin{center}
Trình duyệt $\rightarrow$ Web app (HTTPS) $\rightarrow$ API (HTTPS, JSON, JWT trong header \code{Authorization}) $\rightarrow$ PostgreSQL (kết nối nội bộ)
\end{center}

Web app không truy cập database trực tiếp. Giỏ hàng lưu ở trình duyệt, nhưng giá cuối và kiểm tra nguyên liệu luôn do API thực hiện.

\subsubsection{Các lớp và module của API}
Mỗi module NestJS có ba lớp:
\begin{enumerate}
\item \textbf{Controller:} nhận HTTP request, kiểm tra JWT và vai trò bằng guard, validate DTO bằng \code{ValidationPipe} với \code{whitelist} và \code{forbidNonWhitelisted}.
\item \textbf{Service:} nghiệp vụ, transaction, máy trạng thái đơn.
\item \textbf{Truy cập dữ liệu:} Prisma Client qua \code{PrismaService} dùng chung.
\end{enumerate}

Controller chỉ gọi service của module mình. Service gọi module khác qua interface công khai, không truy cập trực tiếp bảng của module khác ngoài transaction dùng chung.

{\small
\begin{longtable}{@{}L{2.2cm}L{7cm}L{4.7cm}@{}}
\caption{Các module của API và lớp domain tương ứng}\label{tab:module}\\
\toprule
\textbf{Module} & \textbf{Trách nhiệm} & \textbf{Lớp domain (class diagram)} \\ \midrule \endfirsthead
\toprule \textbf{Module} & \textbf{Trách nhiệm} & \textbf{Lớp domain (class diagram)} \\ \midrule \endhead
\bottomrule \endfoot
\code{auth} & Đăng ký, đăng nhập, phát JWT, guard theo vai trò & User, Role \\
\code{catalog} & Sản phẩm, topping, tính số ly còn đặt được theo nguyên liệu & Product, Topping \\
\code{inventory} & Công thức, định lượng theo size/topping/đường/đá, trừ và hoàn kho, nhập kho & Ingredient, RecipeItem, ToppingRecipeItem, OrderIngredientUse \\
\code{orders} & Tạo đơn, máy trạng thái, lịch sử, hủy đơn, tự hủy đơn quá hạn & Order, OrderItem, OrderStatus \\
\code{payments} & Thanh toán Ví/Thẻ/QR, idempotency, adapter cổng thanh toán & Payment, PaymentIntent, PaymentAccount \\
\code{promotions} & Kiểm tra và áp mã khuyến mãi & Promotion \\
\code{loyalty} & Cộng và trừ điểm, đổi quà & LoyaltyTransaction, Reward, RewardRedemption \\
\code{health} & \code{GET /health} kiểm tra API và kết nối database & -- \\
\end{longtable}
}

\textbf{Đối chiếu thiết kế với mã nguồn.} Khi hiện thực, tên và ranh giới module thay đổi nhẹ so với bảng \ref{tab:module}. Thư mục \code{apps/api/src} có các module sau:
\begin{itemize}
\item \code{auth}, \code{products} (thay cho \code{catalog});
\item \code{ingredients} (API kho cho nhân viên), \code{inventory} (các hàm tính nguyên liệu, số ly, trừ kho);
\item \code{orders}, \code{payments} (gồm \code{vnpay});
\item \code{loyalty}, \code{rewards}, \code{health}.
\end{itemize}

Phần khuyến mãi không tách module riêng mà nằm trong \code{orders/\allowbreak{}pricing.\allowbreak{}ts}, vì chỉ dùng khi tính giá đơn. Các thay đổi này không ảnh hưởng quy tắc phụ thuộc. Nhóm ghi lại ở đây để người đọc mã nguồn không bị nhầm.

\subsection{Các cơ chế xuyên suốt}
{\small
\begin{longtable}{@{}L{2.8cm}L{8.4cm}L{2.7cm}@{}}
\caption{Các cơ chế xuyên suốt}\label{tab:xuyensuot}\\
\toprule
\textbf{Cơ chế} & \textbf{Cách làm} & \textbf{Yêu cầu} \\ \midrule \endfirsthead
\toprule \textbf{Cơ chế} & \textbf{Cách làm} & \textbf{Yêu cầu} \\ \midrule \endhead
\bottomrule \endfoot
Transaction và tồn kho & \code{prisma.\allowbreak{}\$transaction}; trừ kho bằng \code{updateMany} có điều kiện \code{stock >= lượng cần}; số dòng cập nhật nhỏ hơn số nguyên liệu thì rollback & FR-12, UC-01 E4 \\
Idempotency & Header \code{Idempotency-Key}; ràng buộc \code{UNIQUE(user\_\allowbreak{}id, idempotency\_\allowbreak{}key)} trên \code{payments}; gửi lại cùng key trả kết quả đã lưu & FR-09 \\
Xác thực, phân quyền & JWT có hạn; mật khẩu và PIN băm bằng bcrypt; \code{RolesGuard} cho STAFF & FR-05, FR-06, NFR-02 \\
Cấu hình & Biến môi trường, kiểm tra khi khởi động; \code{.env.example} trong repo, \code{.env} không commit & NFR-03 \\
Logging & Logger của NestJS, mỗi request ghi method, path, status, thời gian; không ghi mật khẩu, PIN, token & NFR-09 \\
Health check & \code{GET /health} trả \code{\{status: "ok", db: "up"\}} & NFR-09 \\
Lỗi & Exception filter trả JSON thống nhất; frontend chuyển thành thông báo tiếng Việt & NFR-06 \\
Migration & Prisma Migrate; seed tái tạo dữ liệu mẫu & NFR-10 \\
\end{longtable}
}

\subsection{Lý do chọn, ưu điểm, hạn chế và trade-off}
\textbf{Lý do chọn.} Đề bắt buộc Next.js, NestJS và PostgreSQL. Kiến trúc phân lớp với module theo domain là cách tổ chức mặc định của NestJS, phù hợp với nhóm 4 người trong thời gian ngắn.

\textbf{Ưu điểm:}
\begin{itemize}
\item Một API, một database: dễ dùng transaction cho đặt đơn và thanh toán, dễ chạy bằng Docker Compose.
\item Module theo domain khớp class diagram, dễ chia việc và dễ truy vết yêu cầu $\rightarrow$ module $\rightarrow$ test.
\item Cổng thanh toán qua interface: thêm VNPay sandbox mà không sửa nghiệp vụ đơn hàng.
\end{itemize}

\textbf{Hạn chế:}
\begin{itemize}
\item Monolith: mọi module scale cùng nhau; một lỗi nghiêm trọng có thể làm dừng cả API.
\item Frontend và backend là hai dịch vụ, phải cấu hình CORS và URL API theo môi trường.
\end{itemize}

{\small
\begin{longtable}{@{}L{3.8cm}L{4.4cm}L{5.7cm}@{}}
\caption{Trade-off kiến trúc đã chấp nhận}\label{tab:tradeoff}\\
\toprule
\textbf{Lựa chọn} & \textbf{Thay vì} & \textbf{Lý do} \\ \midrule \endfirsthead
\toprule \textbf{Lựa chọn} & \textbf{Thay vì} & \textbf{Lý do} \\ \midrule \endhead
\bottomrule \endfoot
Modular monolith & Microservices & Quy mô nhỏ, cần transaction xuyên module (đơn, kho, thanh toán, điểm) \\
Trừ kho có điều kiện trong transaction & Khóa bi quan (\code{SELECT .\allowbreak{}.\allowbreak{}.\allowbreak{} FOR UPDATE}) hoặc hàng đợi & Đơn giản, đủ an toàn cho mức tải của một cửa hàng, kiểm thử được bằng test đồng thời \\
Idempotency lưu trên bảng \code{payments} & Bảng idempotency riêng & Ít bảng hơn, mỗi lần thanh toán đúng một bản ghi \\
Giỏ hàng lưu ở trình duyệt & Lưu giỏ hàng trên server & Giảm API và bảng; giá cuối vẫn do server tính khi tạo đơn \\
\end{longtable}
}

\subsection{Quyết định kiến trúc (ADR)}\label{sec:adr}
Mỗi ADR có bốn phần: bối cảnh, quyết định, hệ quả và phương án đã cân nhắc. Bảng \ref{tab:adr} tóm tắt tám ADR. Các mục sau trình bày chi tiết ba ADR ảnh hưởng nhiều nhất đến nghiệp vụ và triển khai.

\begin{longtable}{@{}L{1.1cm}L{4.2cm}L{4.4cm}L{4.1cm}@{}}
\caption{Tóm tắt các ADR}\label{tab:adr}\\
\toprule
\textbf{ADR} & \textbf{Quyết định} & \textbf{Hệ quả chính} & \textbf{Phương án bị loại} \\ \midrule \endfirsthead
\toprule \textbf{ADR} & \textbf{Quyết định} & \textbf{Hệ quả chính} & \textbf{Phương án bị loại} \\ \midrule \endhead
\bottomrule \endfoot
001 & Monorepo \code{apps/web}, \code{apps/api}, \code{docs}, \code{tests} với npm workspaces & Một PR thay đổi được API, giao diện và tài liệu; CI chung & Hai repository riêng \\
002 & Backend NestJS phân lớp, module theo domain (modular monolith) & Transaction xuyên module trong một database; scale cùng nhau & Microservices (cần saga) \\
003 & PostgreSQL với Prisma ORM và Prisma Migrate & Kiểu sinh tự động; \code{\$transaction}, \code{updateMany} có điều kiện; migration commit vào repo & TypeORM; SQL thuần \\
004 & Interface \code{PaymentGateway}: cổng mô phỏng và VNPay sandbox & Nghiệp vụ không phụ thuộc cổng; test dùng cổng mô phỏng & Gọi VNPay trực tiếp trong service đơn \\
005 & Trừ kho có điều kiện trong transaction trước khi gọi cổng; idempotency trên bảng \code{payments} & Kho không âm, không bán vượt; giữ transaction trong lúc gọi cổng mô phỏng & \code{SELECT .\allowbreak{}.\allowbreak{}.\allowbreak{} FOR UPDATE}; bảng idempotency riêng \\
006 & JWT có hạn, bcrypt, \code{RolesGuard}; đơn người khác trả 404 & API stateless; không thu hồi JWT trước hạn & Session cookie trên server \\
007 & Docker, GHCR, Render; CI/CD bằng GitHub Actions & Image đã test là image được triển khai; Render Free ngủ khi không có truy cập & Render tự build; VPS tự quản \\
008 & Cấu hình qua biến môi trường, log ra stdout, \code{GET /health} & Một image chạy mọi môi trường & File cấu hình theo môi trường trong repo \\
\end{longtable}

\subsubsection{ADR 004 -- Cổng thanh toán qua interface}
\textbf{Bối cảnh.} Ví, Thẻ, QR dùng cổng mô phỏng (FR-08, FR-18), cần điều khiển được kết quả để kiểm thử. VNPay sandbox là yêu cầu Could (FR-21). Cổng thật nằm ngoài phạm vi.

\textbf{Quyết định.} Định nghĩa interface \code{PaymentGateway} với hai cài đặt:
\begin{itemize}
\item \code{MockPaymentGateway}, chạy trong API:
  \begin{itemize}
  \item Ví và Thẻ trả thành công hoặc thất bại theo dữ liệu thử.
  \item QR dùng \code{PaymentIntent} (token, hạn dùng) và \code{PaymentAccount} thử nghiệm có PIN băm và số dư.
  \end{itemize}
\item \code{VnpayGateway}: tạo URL thanh toán sandbox; ở Return URL kiểm tra chữ ký HMAC-SHA512 và số tiền.
\end{itemize}
Thông tin VNPay chỉ nằm trong biến môi trường. Cổng được gọi \emph{sau khi} đã trừ kho có điều kiện.

\textbf{Hệ quả.} Nghiệp vụ đơn hàng không phụ thuộc cổng cụ thể. Bỏ VNPay vẫn không ảnh hưởng MVP. Hạn chế: cổng mô phỏng không phản ánh hết độ trễ và lỗi mạng của cổng thật.

\subsubsection{ADR 005 -- Tồn kho đồng thời và idempotency}
\textbf{Bối cảnh.} Nhiều khách có thể thanh toán cùng lúc cho cùng nguyên liệu. Không được bán vượt tồn kho, không được trừ tiền khi đã hết nguyên liệu (UC-01 E4), và gửi lại yêu cầu thanh toán không được tạo thêm thanh toán (FR-09).

\textbf{Quyết định:}
\begin{enumerate}
\item \textbf{Tạo đơn:} chỉ kiểm tra đủ nguyên liệu và ghi nhu cầu vào \code{OrderIngredientUse}, chưa trừ kho.
\item \textbf{Thanh toán}, trong một transaction:
  \begin{enumerate}
  \item Trừ kho bằng \code{UPDATE ingredients SET stock = stock - :qty WHERE id = :id AND stock >= :qty} cho từng nguyên liệu.
  \item Nếu có nguyên liệu không cập nhật được: rollback, không gọi cổng, đơn PAYMENT\_FAILED kèm lý do.
  \item Nếu đủ: gọi cổng thanh toán.
    \begin{itemize}
    \item Thành công: ghi \code{Payment}, chuyển PAID, cộng điểm, commit.
    \item Thất bại: rollback (kho hoàn lại), ghi \code{Payment} FAILED, đơn PAYMENT\_FAILED.
    \end{itemize}
  \end{enumerate}
\item \textbf{Idempotency:}
  \begin{itemize}
  \item Client gửi header \code{Idempotency-Key}.
  \item Bảng \code{payments} có \code{UNIQUE(user\_\allowbreak{}id, idempotency\_\allowbreak{}key)} và \code{request\_hash}.
  \item Gặp lại khóa với cùng nội dung: trả kết quả đã lưu kèm header \code{Idempotent-Replayed: true}.
  \item Cùng khóa nhưng khác nội dung: trả 409.
  \end{itemize}
\end{enumerate}

\textbf{Hệ quả.} Tồn kho không âm và số đơn PAID không vượt quá nguyên liệu. Điều này được kiểm chứng bằng test gửi đồng thời nhiều yêu cầu (INT-P07, INT-P08). Với VNPay (redirect, bất đồng bộ), kho được trừ khi xử lý Return URL hoặc IPN hợp lệ, theo cùng điều kiện.

\subsubsection{ADR 007 -- Triển khai bằng Docker, GHCR và Render}
\textbf{Bối cảnh.} Đề yêu cầu:
\begin{itemize}
\item production truy cập qua HTTPS;
\item có container registry;
\item pipeline CD gồm build image, push registry, deploy, smoke test;
\item có cách rollback.
\end{itemize}
Nhóm không có máy chủ riêng.

\textbf{Quyết định:}
\begin{itemize}
\item \textbf{Local:} \code{docker compose up} chạy \code{db}, \code{api}, \code{web}.
\item \textbf{CI:} chạy trên mỗi PR và mỗi lần push lên \code{main}.
\item \textbf{CD:} chạy khi push tag \code{v*}, gồm CI $\rightarrow$ build image \code{api}, \code{web} $\rightarrow$ push GHCR $\rightarrow$ deploy hook Render $\rightarrow$ smoke test.
\item \textbf{Production:} hai Web Service trên Render và một Render PostgreSQL chỉ nhận kết nối nội bộ.
\item \textbf{Secret:} chỉ nằm trong biến môi trường của Render và GitHub Secrets.
\item \textbf{Rollback:} triển khai lại image của tag cũ trên GHCR. Migration chỉ thêm, không xóa cột trong cùng phiên bản, để rollback image an toàn.
\end{itemize}

\textbf{Hệ quả.} Image đã test là image được triển khai, và truy vết được phiên bản qua tag. Hạn chế: gói Free của Render ngủ khi không có truy cập, lần gọi đầu chậm.

\subsection{Task 3.2 -- Deployment diagram}
Deployment diagram (hình \ref{fig:deployment}) thể hiện đúng hạ tầng thật của bản production:
\begin{itemize}
\item GitHub repository và GitHub Actions (CI 2 job, CD theo tag);
\item GitHub Container Registry chứa image \code{brewlite-as2-api}, \code{brewlite-as2-web};
\item Render (Singapore) gồm Web Service web, Web Service API và Render PostgreSQL 16 chỉ kết nối nội bộ;
\item VNPay sandbox bên ngoài;
\item thiết bị người dùng (điện thoại, máy tính) chạy trình duyệt hoặc PWA, truy cập qua HTTPS cổng 443 vào lớp edge (TLS) của Render.
\end{itemize}

Môi trường local bằng Docker Compose được thể hiện riêng.

\hinh{deployment-diagram.pdf}{0.75\linewidth}{Deployment diagram (\code{docs/\allowbreak{}architecture/\allowbreak{}deployment-diagram.\allowbreak{}puml})}{fig:deployment}

Thông tin chi tiết của production (tên dịch vụ, biến môi trường, cách bảo mật database) ở mục \ref{sec:production}.
